Hosted image editors often modify facial regions outside the anatomy named in a
cosmetic-surgery preview request, and every training-free localization method in
the literature needs model internals that a hosted endpoint does not expose. We
conduct a pilot benchmark of the control that remains available on the client:
six hosted editing configurations and one mask-specific inpainting endpoint on
facelift-style jaw--neck and rhinoplasty edits, under a three-rung ladder of
prompt only, a landmark-derived client-side masked composite, and
endpoint-native masked inpainting where available. Across \NumAttemptedEdits{}
attempted cells, \NumScoredEdits{} outputs were scoreable. ArcFace cosine
measures identity retention, while a CIELAB ratio measures change in the target
region relative to a predefined facial keep zone. Among the
\NumLocalizationFaces{} frontal faces for which regional scoring was available,
compositing increased localization over the paired prompt-only output by a
median of \MedPairedLocGain{} (95\% face-clustered bootstrap interval
\PairedLocGainLow--\PairedLocGainHigh) while retaining comparable on-target
pixel change. Model choice produced a visible trade-off between edit strength
and identity retention, and the single tested inpainting endpoint did not
outperform post-hoc compositing. The claims are scoped to control rather than
clinical accuracy: no surgeon ratings were collected, and each cell contains
one generated output. Within that scope, region
confinement for surgical previews requires no model access; it can be enforced
on the client alone, uniformly across the tested editors and at low provider
cost.
